\section{Multi-domain Data Pipeline}
\subsection{Long-horizon Search}
\label{sec:search_data}
We instantiate the knowledge-action graph (KAG) in the search domain by constructing search question data over a large wiki database. Here $\mathcal{C}_d$ is a corpus of wiki pages and paragraph-level evidence, $\mathcal{A}_d$ consists of entity transitions induced by hyperlinks and controlled graph walks, $\mathcal{O}_d$ contains the retrieved page text and connecting evidence, and $\mathcal{V}_d$ checks whether the answer entity and its supporting path are recoverable from the provided context. This instantiation targets the information acquisition and evidence verification abilities defined in Section~\ref{sec:domain-infrastructure}. The pipeline first converts wiki entries into a directed corpus graph, where each entry is an entity node and inter-entry references define edges. For each node, the pipeline retrieves page text, paragraph structure, outgoing anchors, canonical titles, and optional structural statistics such as in-degree, out-degree, and text length. Non-content tail sections and anchors from list-style paragraphs are filtered to reduce noisy transitions, so the resulting graph can provide both the search space and the provenance records required by $\mathcal{G}_d$.

\textbf{Relation-chain generation.} Relation chains are generated through controlled random walks over the corpus graph and the corresponding domain KAG. Candidate next nodes are selected from outgoing anchors after removing visited entities, near-duplicate title variants, disambiguation pages, pages without valid text, and pages with insufficient outgoing links. Degree and text-length constraints further control node quality and graph diversity. An LLM selector chooses one entity from the remaining candidates, balancing local path coherence with optional cross-domain topic shifts. If a walk reaches a dead end, a short auxiliary walk recovers a qualified continuation node. Each accepted transition is stored as an action-like record linking the current entity state, the selected hyperlink action, the observed target page, and the paragraph evidence that justifies the transition.

\textbf{Question-answer pair generation.} Each completed chain is serialized with its ordered entity list, node texts, and paragraph-level evidence for adjacent entities. The generator treats the final entity as the answer, rewrites the chain into a coherent natural-language question, masks proper nouns and other identifying information, and requires the solver to recover the masked answer entity. A verifier then checks the answer identity and evidence support against the serialized chain. Each chain, therefore, yields a verifiable search-domain training instance that requires following indirect long-context evidence rather than relying on direct name matching, while preserving the provenance, action path, observation text, and verifier target needed by the KAG.

\textbf{Trajectory collection and quality control.} Search trajectories are collected by allowing strong models to execute deep-research tasks with the \texttt{search}, \texttt{read\_page}, and \texttt{code} tools in the real Internet environment. The \texttt{search} tool queries a commercial search engine and returns the top results per query. The \texttt{read\_page} tool extracts web-page content and uses an LLM to summarize the extracted information before it is returned to the agent context. The \texttt{code} tool allows the model to write and execute Python code in a sandboxed environment, supporting intermediate computation and data processing during the search process. The maximum context window is set to 256K tokens, and no additional constraint is imposed on the number of tool calls within a single turn. During post-processing, we remove trajectories with wrong answers, overly short interaction histories, or obvious guessing behavior. For trajectories that reach the maximum turn limit, we roll back to the previous turn and use an explicit user prompt to force the model to produce an answer. The retained trajectories provide long-horizon supervision for search behavior that combines query formulation, page reading, evidence integration, and answer verification under realistic web observations.

%%%%%%%%%%%%%%%%%%%%%%%%%%%%%%%%%

\subsection{Machine Learning Engineering}
\label{sec:mle_data_pipeline}
Following the knowledge-action infrastructure in Section~\ref{subsec:knowledge-action-supervision}, we instantiate machine learning engineering (MLE) as a KAG over executable solution search. Here $\mathcal{C}_d$ contains competition specifications, datasets, and execution environments; $\mathcal{A}_d$ captures solution edits, experiment runs, tree navigation, node invalidation, and answer commitment; $\mathcal{O}_d$ records logs, metrics, artifacts, and generated submissions; and $\mathcal{V}_d$ provides grader scores, submission-format checks, and metric-reliability checks. Because a candidate solution is executable code whose quality is known only after grading, each trajectory becomes a verifier-guided expansion of the solution graph.

\textbf{Task sources.}
We assemble gradeable competitions from MLE-Dojo~\citep{mledojo} and ended Kaggle competitions processed by our automated framework. MLE-Dojo provides curated Kaggle-style tasks with held-out answers and local graders across tabular, vision, NLP, audio, and time-series settings. For ended competitions, our framework curates the task description and leaderboard information, re-splits public data into fresh train/test partitions, constructs private answers, and synthesizes a local evaluator with submission validation. This produces a diverse and refreshable pool of optimization tasks while reducing overfitting to any static benchmark distribution.

\textbf{Agentic harness.}
Trajectories are generated in an agentic harness that grows a tree of executable solution nodes, following prior solution-search workflows such as MLEvolve~\citep{mlevolve}. Writing a full script opens a new root, patching a node spawns a child, and executing a node attaches observations such as logs, exceptions, metrics, artifacts, and submission validity. This tree forms the executable core of the MLE-domain KAG: branches encode alternative attempts, verifier outcomes guide expansion, and failed or invalidated nodes provide negative evidence for later decisions.

The harness exposes this process through a compact tool interface, summarized in Table~\ref{tab:mle_tools}, covering code authoring, execution, tree navigation, answer management, persistent memory, and delegated analysis. For long runs, an isolated \texttt{analyze} sub-agent investigates data or results and returns a report, while context compaction summarizes earlier steps into a digest.

\textbf{Trajectory collection and quality control.}
We collect teacher trajectories over the task pool with multiple seeds and prompt variants, replay them with the local evaluator, and retain runs that produce valid and competitive committed submissions. After trimming regressive or degenerate segments, we deduplicate the remaining runs and serialize them into a unified message schema with loss masks for teacher-generated content. The serialization preserves node relations, execution observations, verifier outcomes, and committed-answer history, making the solution-search structure recoverable for training.

\begin{table}[ht]
\centering
\caption{\textbf{Our designed tool interface used by \ProjectName, according to the MLE agentic harness.} The tools define a compact action space through which the model performs code authoring, execution, search-tree management, persistent memory, and delegated analysis.}
\label{tab:mle_tools}
\begin{tabular}{@{}l p{0.66\linewidth}@{}}
\toprule
\textbf{Tool} & \textbf{Function} \\
\midrule
\multicolumn{2}{@{}l}{\textit{Code authoring \& execution}}\\
\texttt{write\_full\_code} & Author a complete training script from scratch; opens a new root node (a fresh line of attack). \\
\texttt{patch\_code} & Apply a localized edit to a node's code; spawns a child node, preserving tree history for incremental refinement. \\
\texttt{execute\_code} & Run a node, capture stdout and exceptions, extract its validation metric, and check the emitted submission for validity. \\
\texttt{execute\_bash} & Run a guarded shell command for environment setup and inspection (installs, GPU checks, file operations). \\
\addlinespace
\multicolumn{2}{@{}l}{\textit{Search-tree navigation \& answer management}}\\
\texttt{list\_nodes} & Survey the solution tree: the selected answer, the recent answer trail, invalidated history, and a metric-ranked listing. \\
\texttt{select\_node} & Inspect one node in full (code, plan, output, metric, parent chain) before revisiting or branching from it. \\
\texttt{invalidate\_node} & Exclude a node whose metric is untrustworthy (leakage, overfitting) from ranking and submission. \\
\texttt{update\_answer} & Commit a node as the current submission candidate, written to the canonical path the grader reads. \\
\texttt{get\_current\_answer} & Report the node currently committed as the answer. \\
\addlinespace
\multicolumn{2}{@{}l}{\textit{Persistent memory}}\\
\texttt{write\_notes} / \texttt{read\_notes} & Append to / re-read a notebook that survives context compaction (decisions, failed strategies and why, hypotheses). \\
\addlinespace
\multicolumn{2}{@{}l}{\textit{Sub-agent}}\\
\texttt{analyze} & Spawn an isolated analysis sub-agent that explores data and results in its own context window and returns a single structured report. \\
\bottomrule
\end{tabular}
\end{table}

%%%%%%%%%%%%%%%%%%%%%%%%%%%%%%%%%

\subsection{Scientific Reasoning and Research}
\label{sec:science-data-pipeline}
Following the knowledge-action infrastructure in Section~\ref{subsec:knowledge-action-supervision}, we instantiate scientific reasoning as a domain-specific KAG over scientific problem-solving processes. Formally, given the scientific knowledge-action graph $\mathcal{G}_d = (\mathcal{C}_d, \mathcal{A}_d, \mathcal{O}_d, \mathcal{V}_d)$,  $\mathcal{C}_d$ contains scientific problem components including problem statements, topics, keywords, and solution structures; $\mathcal{A}_d$ consists of reasoning actions such as decomposition, transformation, retrieval, and computation; $\mathcal{O}_d$ captures intermediate reasoning states, external evidence, and execution outputs; and $\mathcal{V}_d$ contains verification signals over correctness, consistency, and scientific validity.

\textbf{Problem Construction.} We first collect a large-scale pool of scientific problems spanning fundamental domains such as mathematics, physics, and related areas. Based on this seed pool, we construct a knowledge-action graph, where each seed is represented as interconnected nodes encoding its problem statement, keywords, and solution components, forming an initial knowledge-aligned subgraph. Built upon this graph, we perform self-evolving scientific KAG enhancement introduced in Section \ref{subsec:knowledge-action-supervision} via self-play graph search and expansion, where local subgraphs are sampled and systematically rewritten to improve scientific problems along two complementary directions: (i) \emph{harder reasoning variants}, which increase required domain knowledge depth, introduce complex symbolic structures, and extend multi-step derivations; and (ii) \emph{interaction-enriched variants}, which inject cross-domain knowledge, incorporate concepts requiring external retrieval, and increase computational demands such as code-based numerical execution. For problems that already exhibit strong reasoning ability or strong interaction requirements, we retain their original form without modification. Through this graph-driven expansion process, we construct an enhanced scientific problem corpus of approximately 15K instances, characterized by higher difficulty and stronger interaction, thereby providing a data foundation for long-horizon scientific reasoning.

\textbf{Trajectory Generation.} Based on the enhanced scientific problem corpus, we construct both no-tool and tool-augmented reasoning trajectories as training data for long-horizon scientific reasoning. For no-tool reasoning, we distill high-quality pure reasoning trajectories from a strong model, consisting of multi-step derivations, symbolic transformations, and final answers, and retain only those verified as correct to improve pure reasoning capability. For tool-augmented reasoning, we build an interactive reasoning framework equipped with four external tools: \texttt{search}, \texttt{visit}, \texttt{code}, and \texttt{scholar}. The \texttt{search} tool retrieves relevant web information for factual grounding; the \texttt{visit} tool enables inspection of specific web pages or documents; the \texttt{code} tool supports numerical computation, symbolic calculation, equation solving, and simulation; and the \texttt{scholar} tool provides access to academic literature. Using this framework, we distill tool-augmented trajectories from a strong model, recording intermediate reasoning steps, tool calls, observations, computations, and final answers, and filter them by final-answer correctness to ensure high-quality interaction data. Overall, no-tool and tool-augmented trajectories provide complementary supervision for strengthening reasoning and interaction capabilities in long-horizon scientific problem solving.

%%%%%%%%%%%%%%%%%%%%%%%%%%%%%%%%%
\subsection{Instruction Following}
\label{sec:if_data}

Following the knowledge-action infrastructure in Section~\ref{subsec:knowledge-action-supervision}, we instantiate instruction following as a KAG-style supervision problem over constraints, long-context evidence, and locally introduced rules. Here $\mathcal{C}_d$ contains user instructions, verifiable constraints, long-document chunks, document-level entities and facts, injected in-context rules, distractors, and answer candidates; $\mathcal{A}_d$ consists of KAG-level operations such as constraint parsing, evidence selection, local-rule application, distractor rejection, and final answer formatting; $\mathcal{O}_d$ contains selected evidence, parsed constraint states, injected rule states, and candidate answer states; and $\mathcal{V}_d$ contains automatic validators for constraint satisfaction, answer matching, evidence dependence, and consistency with injected in-context information. This instantiation targets constraint tracking, evidence verification, and long-context adaptation.

\textbf{Task Construction.}
We construct the instruction-following corpus from two complementary sources with the corresponding domain KAG. The first source is a high-quality subset of 13K multi-constraint instruction samples from NVIDIA's Nemotron-RL instruction-following dataset~\citep{nemo-gym}. These examples are built from prompts in WildChat-1M~\citep{zhao2024wildchat} and instructions from the Open-Instruct codebase, and cover automatically verifiable constraints such as length, format, keywords, language, punctuation, and paragraph structure. In the KAG view, each constraint is represented as a condition node linked to the required output behavior and its validator, so the resulting data directly supervises whether the model can parse and satisfy explicit user requirements.

The second source is a self-built long-context learning pipeline that produces 10K verified long-context QA instances. We first structurally parse long documents and extract entities, attributes, relations, numerical values, and other salient facts to construct document-level evidence graphs compatible with the KAG representation. Based on these graphs, we synthesize multi-hop QA tasks that require combining evidence across multiple factual nodes. We then inject local in-context rules or distractors, such as temporary protocols that override document rules, misleading framing statements, or constraints that conflict with prior knowledge. These injected signals are linked to the relevant evidence nodes, rule nodes, and verifier checks. Candidate tasks are converted into a unified multiple-choice QA format and filtered by automatic validation to ensure that the correct answer depends on both the long-context evidence chain and the injected in-context information. This yields verifiable supervision for locating dispersed evidence, applying local rules, and resisting unsupported distractors.

%%%%%%%%%%%%%%%%%%%%%%%%%%%%%%%%%
\subsection{Tool Calling}
\label{sec:tool-calling}

Following the knowledge-action infrastructure in Section~\ref{subsec:knowledge-action-supervision}, we instantiate tool calling as a KAG over executable interactions. Here $\mathcal{C}_d$ contains tool schemas, environment states, optional simulated user profiles, and task resources; $\mathcal{A}_d$ consists of schema-grounded calls and clarification actions; $\mathcal{O}_d$ contains tool returns, state updates, errors, and user feedback; and $\mathcal{V}_d$ contains available schema, state, grounding, and goal-completion checks. Each task can induce multiple candidate long-horizon state--action--observation chains whose later decisions depend on earlier observations.

\textbf{Task Construction.} The task pool and tool interface are constructed jointly through tool extraction, tool-interaction graph construction, graph-compositional task synthesis, and solvability assessment. We extract candidate interfaces from scientific, web, repository, database, and locally simulated tool-usage settings~\citep{barres2025tau2,he2025vitabench}. The graph connects tools, resource/state types, observation fields, and verifier targets, with edges denoting executable dependencies such as schema compatibility, precondition--effect relations, shared resources, state transitions, or support for a completion check. Task synthesis is formulated as constrained graph search over this dependency graph: the generator selects connected tool subgraphs or dependency paths anchored by target states, answers, or verifier conditions, and then renders them into user instructions. This avoids arbitrary tool concatenation: later calls must rely on objects, constraints, or observations produced earlier. Candidate tasks are screened for argument availability, initial-state reachability, sandbox executability, and verifier coverage. Underspecified, prompt-only, schema-incompatible, or unverifiable tasks are revised or discarded. Each retained task stores the user goal, selected tool subgraph, initial state, hidden support path, and completion criterion or verifier.

\textbf{Trajectory Generation.} Trajectory generation is performed in a Tool Sandbox, which exposes the selected tool subgraph and maintains the evolving environment state. Solver backends explore the candidate space over multiple turns by choosing tools, binding arguments, interpreting observations, and deciding whether to continue, ask for clarification, or produce the final response. When preferences, missing constraints, or task-bounded clarifications are required, a simulated user participates in the loop. For the same user goal and verifier target, including the same question--answer pair in QA-style tasks, we generate multiple candidate trajectories through different tool choices, actions, or clarification turns. This process is treated as verifier-guided graph search over the trajectory space. Available verifier or judge modules score or rank these candidates by call format, schema typing, state consistency, observation grounding, constraint satisfaction, and goal completion. Invalid schemas, unsupported arguments, unrecovered failures, or ungrounded responses are rejected or routed to repair and resampling. Accepted runs are kept as KAG-aligned message trajectories containing assistant turns, tool calls, observations, user feedback, state changes, and verifier outcomes.